\begin{example}\label{ex:n20}
For $n=20$, $F_{20}=6765=3\cdot 5\cdot 11\cdot 41$. The relevant atomic
factors are $3\in\Atom(4)$, $5\in\Atom(5)$, $11\in\Atom(10)$,
$41\in\Atom(20)$. The witness covers are $\{41\}$, $\{3,5\}$, $\{3,11\}$,
yielding $\mathcal M_{20}=\{15,33,41\}$ and $A(20)=11$.
Indeed, with coordinates $(2,5)$, the ranks $20$, $4,5$, and $4,10$
respectively cover both coordinates and have the singleton essential supports
required by Definition~\ref{article-def:witness-cover}; the value $A(20)=11$ is the
exhaustive divisor/rank count reported in Table~\ref{tab:birth-layer-data-generated}.
\end{example}

\begin{proposition}[Exact support-three realization at $n=30$]
\label{prop:n30-five-type-classification}
Let the prime coordinates of $30$ be labelled by $(2,3,5)$. Then
\[
 \mathcal M_{30}=\{20,22,31,244,671\}.
\]
These five generators realize, respectively, the five unlabelled types
$\Gamma_8,\Gamma_7,\Gamma_1,\Gamma_5,\Gamma_4$.  No generator in
$\mathcal M_{30}$ has type $\Gamma_3$, $\Gamma_6$, or $\Gamma_9$.
Consequently exactly five, not all eight, of the Fibonacci-admissible
three-coordinate types occur at $n=30$.
\end{proposition}

\begin{proof}
The exact factorization
\[
 F_{30}=832040=2^3\cdot5\cdot11\cdot31\cdot61
\]
reduces the atomic inventory to the prime powers
$2,4,8,5,11,31,61$. Proposition~\ref{prop:prime-power-ladders} gives the
ranks of $2,4,8$, and $5$. For the remaining primes,
Proposition~\ref{article-prop:rank-basics} reduces the rank computation to the proper
divisors of $10$, $15$, and $30$: one has
\[
 F_{10}=55,\qquad F_{15}=610,\qquad F_{30}=832040,
\]
and the required nondivisibilities are
\[
\begin{aligned}
11&\nmid F_d &&(d\in\{1,2,5\}),&
61&\nmid F_d &&(d\in\{1,3,5\}),\\
31&\nmid F_d &&(d\in\{1,2,3,5\}),&
31&\nmid F_d &&(d\in\{6,10,15\}).
\end{aligned}
\]
Hence
\[
\begin{array}{c|rrrrrrr}
\theta&2&4&8&5&11&31&61\\ \hline
\alpha(\theta)&3&6&6&5&10&30&15.
\end{array}
\]
Thus $8$ is not atomic, because $\alpha(8)=\alpha(4)$, while the other six
entries are atomic.  Their nonempty slots, with coordinate names used as
subscripts, are exactly
\[
 P_{235}=\{31\},\quad P_{25}=\{11\},\quad P_{35}=\{61\},
 \quad S_2=\varnothing,\quad S_3=\{2\},\quad S_5=\{5\},
 \quad L_{2,23}=\{4\}.
\]
Indeed, a prime atom has essential support equal to its full support, which
gives the three displayed exact-rank prime slots and the singleton atoms $2$
and $5$. The only nonprime atom is $4$; lowering it to $2$ changes its rank
from $6$ to $3$, so its full support is $\{2,3\}$ and its essential support
is $\{2\}$.  This also proves that every slot not displayed above is empty.

Applying the if-and-only-if slot test of
Lemma~\ref{lem:three-coordinate-nonemptiness}, the nonempty row fillings are
precisely
\[
\begin{array}{c|c|c}
\text{type}&\text{cover}&\text{product}\\ \hline
\Gamma_1&\{31\}&31\\
\Gamma_4&\{11,61\}&671\\
\Gamma_5&\{4,61\}&244\\
\Gamma_7&\{2,11\}&22\\
\Gamma_8&\{5,4\}&20.
\end{array}
\]
All selected underlying primes are distinct, so every filling passes the
coprimeness condition.  Conversely, the slot test is exhaustive.  More
explicitly, $\Gamma_3$ requires a full-support ladder slot $L_{i,235}$, and
none exists; $\Gamma_6$ requires two nonempty proper ladder slots with
distinct essential coordinates, whereas $L_{2,23}$ is the only proper ladder
slot; and $\Gamma_9$ requires all three singleton slots, whereas
$S_2=\varnothing$.  Hence those three types are impossible at $n=30$, and
the five products displayed above exhaust $\mathcal M_{30}$.
\end{proof}

\subsection*{Substructures of nonexceptional birth layers}

\begin{corollary}\label{cor:canonical-principal-interval}
Let $n\ge3$, $n\notin\{6,12\}$, and choose an exact-rank prime $\pi_n$ with
$\alpha(\pi_n)=n$. Then the principal interval
\[
[\pi_n,F_n]_{D_n}:=\{q\in D_n:\ \pi_n\mid q\}
\]
is contained in $B_n$. In particular $B_n$ contains a divisor-lattice
interval naturally isomorphic to $\operatorname{Div}(F_n/\pi_n)$. This
interval is natural after choosing $\pi_n$; no choice-free canonical interval
is asserted.
\end{corollary}

\begin{proof}
By Theorem~\ref{article-thm:primitive-divisors}, there is a prime
$\pi_n$ with $\alpha(\pi_n)=n$. Thus $\pi_n\in B_n$. Since
Lemma~\ref{article-thm:upper-fiber} says that $B_n$ is an upper set in $D_n$,
every divisor $q$ of $F_n$ with $\pi_n\mid q$ also lies in $B_n$. Division by
$\pi_n$ gives an order-preserving bijection from this interval to
$\operatorname{Div}(F_n/\pi_n)$, with inverse $d\mapsto \pi_n d$.
\end{proof}

\begin{corollary}\label{cor:boolean-skeleton}
Let $n\ge3$, $n\notin\{6,12\}$, and put
\[
\mathcal D_n^*:=\{d:\ d\mid n,\ 3\le d<n,\ d\notin\{6,12\}\}.
\]
For each $d\in\mathcal D_n^*$ choose an exact-rank prime
$\pi_d$ with $\alpha(\pi_d)=d$, and choose an exact-rank prime
$\pi_n$ with $\alpha(\pi_n)=n$. Then
\[
S\longmapsto \pi_n\prod_{d\in S}\pi_d
\qquad(S\subseteq\mathcal D_n^*)
\]
embeds the Boolean lattice $2^{\mathcal D_n^*}$ into $B_n$. Hence
$B_n$ contains a Boolean subposet of rank
\[
|\mathcal D_n^*|\ge \tau(n)-5.
\]
\end{corollary}

\begin{proof}
The exact-rank primes exist by Theorem~\ref{article-thm:primitive-divisors}. They
are distinct, because distinct ranks of apparition cannot belong to the
same prime. For every $d\mid n$, $\pi_d\mid F_d\mid F_n$, so the displayed
products divide $F_n$. They all also divide a multiple of $\pi_n$ inside
$D_n$, and Corollary~\ref{cor:canonical-principal-interval} puts every such
multiple in $B_n$.

Since the primes $\pi_d$ are distinct, divisibility among the displayed
products is exactly inclusion among the subsets $S$. Thus the map is a
Boolean-lattice embedding. Among the $\tau(n)$ divisors of $n$, the only
divisors excluded from $\mathcal D_n^*$ are contained in
$\{1,2,6,12,n\}$, so $|\mathcal D_n^*|\ge\tau(n)-5$.
\end{proof}

%% ====================================================================
